\section{Widening two: when the condition stops answering yes or no}
\label{sec:graded}

\subsection{The move}

A guard returns a truth value, and the communication fires or does not. That is
one choice, not the only one. The second widening is to let a condition return a
value in some other algebra --- a probability, a rate, a cost, a complex
amplitude --- and to interpret that value as \emph{how much} this candidate
happens rather than \emph{whether} it happens.

Before anything else, the reassurance a developer needs:

\begin{proposition}[Conservativity]
\label{prop:cons}
When the value algebra is the Booleans, a graded \texttt{where} clause is
exactly the \texttt{where} clause of \S\ref{sec:cfl}. Nothing about existing
programs changes, and no existing program acquires a new meaning.
\end{proposition}

So this is an extension in the strict sense. A developer who never grades a
condition never encounters any of what follows. A developer who grades one gets
something that turns out to be much larger than a weighted choice.

\subsection{What a value algebra has to supply}

Not every set of values will do. The condition language has conjunction and
disjunction in it, so whatever the values are, there must be an operation to
interpret each. The requirement is mild: a commutative semiring~\cite{GradedWhere}. Multiplication
$\tensorv$ interprets conjunction, addition $\sumv$ interprets disjunction, and
the units behave as true and false do.

\begin{center}
\small
\begin{tabular}{@{}L{0.16\textwidth}L{0.20\textwidth}L{0.26\textwidth}L{0.30\textwidth}@{}}
\toprule
\textbf{Values} & \textbf{$\tensorv$, $\sumv$} & \textbf{A condition means} & \textbf{Running the program means} \\
\midrule
$\mathbb{B}$ &
and, or &
a guard &
executing, as today \\[3pt]
$\Rnn$ &
$\times$, $+$ &
a rate or propensity &
a stochastic simulation \\[3pt]
$[0,1]$ with PLN values &
$\times$, $+$ on the power $\str\conf$ &
a graded belief &
inference driving execution \\[3pt]
Viterbi $(\max,\times)$ &
$\max$, $\times$ &
the best explanation &
finding the most likely path \\[3pt]
tropical $(\min,+)$ &
$\min$, $+$ &
a cost &
finding the cheapest path \\[3pt]
$\mathbb{C}$ &
$\times$, $+$ &
an amplitude &
a quantum simulation \\
\bottomrule
\end{tabular}
\end{center}

Which value algebra a program is read in is not merely a choice of numerics: it
fixes how much of the non-determinism a resolver has to decide, and that
question has a precise
hierarchy behind it~\cite{ChoiceFormulae}. One casualty is worth naming
immediately, because it surprises people. Negation
does not survive. Over $[0,1]$ the map $x \mapsto 1-x$ behaves well enough, but
over the complex numbers there is no canonical negation and De~Morgan fails. The
practical consequence is that the design keeps two sorts of formula: a
\emph{crisp} sort with the full Boolean structure, used for gating, and a
\emph{graded} sort with a semiring only, used for weighting. A \texttt{where}
clause keeps its old job and gains a new slot.

\subsection{The guard becomes a weight}

Recall candidates from \S\ref{sec:transaction}: for a comprehension with data
available, each way of assigning data to patterns is a candidate. In the crisp
setting the condition is evaluated per candidate and admits or refuses it. In
the graded setting the condition is evaluated per candidate and \emph{scores}
it.

That is the entire mechanism, and three worthwhile things fall out of it without
further work.

\paragraph{The propensity is the formula's value, and formulas move.} A crisp
guard is a fixed policy: it fires on the same specimens forever, because the
formula does not change. A graded guard's behaviour is the formula's
\emph{value}, evaluated against the current state --- and the current state
includes whatever the program has learned. Nothing needs to be scheduled.

\paragraph{Weight tables are unnecessary.} A common design for stochastic
execution~\cite{WeightedGSLT} attaches a table of rates to each rule, keyed by conditions, and then
requires those conditions to partition the space so that lookup is
single-valued. A graded clause needs none of that: a formula has one value, so
single-valuedness is free, and overlapping conditions are simply fine.

\paragraph{Plasticity is free.} Because the clause is evaluated against the
current term, weights that change during execution require no update mechanism.
The weight is not stored; it is computed.

\subsection{Three readings of one program}

Here is where the section earns its title. The same program text, under three
different value algebras, is three different kinds of artefact.

\paragraph{Booleans: a state transition.} The program is what it has been all
along --- code running in a node, manipulating state, committing transactions.
Programming in the small.

\paragraph{Non-negative reals: a stochastic simulation.} Each candidate carries a
rate. The total rate at a configuration determines how long the system waits and
which step it takes; this is the standard Gillespie construction~\cite{Gillespie77}, and it turns
the program into a generator of trajectories --- the same move that took the
$\pi$-calculus to the stochastic $\pi$-calculus and to
SPiM~\cite{Priami95,PhillipsCardelli07}, and that the weighted reading of a
language definition generalises~\cite{WeightedGSLT}. Running it many times gives you
the distribution of outcomes rather than one outcome. Programming in the large:
the subject is no longer this program's state but the behaviour of a
\emph{population} of them.

\paragraph{Complex numbers: a quantum simulation.} Each candidate carries an
amplitude~\cite{GWQuantum}, and candidates that reach the same outcome have
their amplitudes added before magnitudes are taken --- which is to say, they interfere. Two
branches that are individually possible can cancel --- the effect Shor's
algorithm is built out of~\cite{Shor97}. A developer's summary:
disjunction stops being monotone, and a false disjunction with true disjuncts
becomes expressible.

\begin{remark}[The simulator is not the interpreter]
\label{sec:simnotinterp}
This matters operationally and is easy to get wrong. The graded readings do not
run inside the node. They are a separate execution path with its own state,
running offline, whose output is a graph or an ensemble rather than a committed
state change. That separation is what makes things affordable there that are not
affordable in the node --- exhaustive construction of the reachable graph, model
checking, fixed points --- and it is what keeps a consensus-critical execution
path free of any of it.
\end{remark}

\subsection{Programming in the small, programming in the large}

The reason to care is not that a simulator exists. Simulators exist. The reason
to care is that \emph{it is the same program text}.

The usual arrangement in industry is two artefacts: the system, written in one
language by one team, and the model of the system, written in another language
by another team, for capacity planning or risk or verification. The model is
always a little wrong, it goes stale immediately, and the discrepancy is
discovered in production.

Here the model and the system are the same source. A developer writes a
comprehension with a graded condition; under the node's resolver it is a
transaction, and under the simulator's resolver it is a model of what a
population of such transactions would do. There is no translation step in which
fidelity is lost, because there is no translation step.

The practical questions this makes askable, in the language the system is
written in:

\begin{itemize}[leftmargin=1.6em,itemsep=2pt]
\item If this contract is deployed a thousand times against that environment,
  what is the distribution of outcomes?
\item Which of these two policies survives longer under a shift in the
  environment?
\item Is there a configuration reachable from here that we would consider a
  failure, and how much would getting there cost an adversary?
\end{itemize}

\subsection{The forager, graded}

Now the running example, all the way through. The forager's clause becomes:

\begin{lstlisting}
new bel, prey, wall in {
  bel!( (1, 0.5) ) |                    // (n, s): one unit of prior evidence

  for( @(n,s) <- bel ; @spec <- prey  where  Chi(spec) * Pow(n,s) ) {
    Break!(spec) | bel!( revise(n, s, spec) )
  } |
  for( @(n,s) <- bel ; @spec <- prey  where  g ) {
    Graze!(*wall) | bel!( (n, s) )
  }
}
\end{lstlisting}

\texttt{Chi(spec)} is the crisp structural predicate of \S\ref{sec:logic},
entering the graded sort as an indicator. \texttt{Pow(n,s)} is a graded ground
formula whose value is the PLN power $\str \cdot \conf$, where $\conf = n/(n+k)$
is confidence and $k$ is the learner's scepticism~\cite{PLN}. \texttt{*} is graded
conjunction. \texttt{g} is a scalar: the value of grazing.

Two comprehensions, each with a weight. The forager opens the specimen with
propensity $\str\conf$ against the wall's $g$. Everything below was measured by
running exactly this.

\paragraph{The bootstrap trap, and why exploration is a disjunct.} A belief with
no evidence has $\conf = 0$, hence power $0$, hence a clause value of zero, hence
a receipt that never fires --- so it never acquires evidence. Measured: a forager
initialised with no prior evidence opens $0.00$ specimens across $400$
encounters and ends with exactly the grazing baseline, against $3543.7$ for the
same forager given one unit of prior evidence.

The remedy is written in the clause language, not in a framework. Either the
belief carries prior evidence $n_0 \geq 1$, or the clause is
$\mathrm{Pow}(n,s) \sumv \varepsilon$ for a small scalar $\varepsilon$. That
second form is $\varepsilon$-greedy exploration, and it is a \emph{disjunct}.

\paragraph{Confidence buys variance, not yield.} Four foragers, two worlds,
$20{,}000$ lives each. In the benign world the crisp forager --- which acts on a
satisfied guard from its first encounter --- has the \emph{highest} mean terminal
holdings, $4084.7$, and the \emph{lowest} survival, $0.8501$. The graded forager
has survival $0.9485$ and holdings $3543.7$. That is the trade a mortal program
actually faces, and a Boolean guard cannot express it.

\paragraph{A Boolean guard has no coordinate in which to record disappointment.}
In the hostile world, where opening is a net loss, the crisp forager survives
$0.29\%$ of the time. It cannot do otherwise: its policy \emph{is} the formula,
and the formula does not change. The graded forager's policy is the formula's
value, which moves:

\begin{center}
\small
\begin{tabular}{@{}lrrrr@{}}
\toprule
opening & $\str$ & $\conf$ & power & $\Pr[\text{open}]$ \\
\midrule
1 & $0.7500$ & $0.3333$ & $0.2500$ & $0.3333$ \\
2 & $0.5000$ & $0.4286$ & $0.2143$ & $0.3000$ \\
3 & $0.3750$ & $0.5000$ & $0.1875$ & $0.2727$ \\
5 & $0.2500$ & $0.6000$ & $0.1500$ & $0.2308$ \\
\bottomrule
\end{tabular}
\end{center}

Strength falls as decoys accumulate, confidence rises, the product falls, and
with it the willingness to open. The forager learns to stop. Note that nothing in
the clause is a schedule: early caution and late decisiveness are consequences of
a coordinate of the truth value, not of a hyperparameter someone tuned.

\paragraph{Your \texttt{where} clause has become a learning rule.} This is the
point at which the exposition arrives somewhere a developer would not have
predicted from \S\ref{sec:cfl}. In a one-dimensional grading, the revision step
that updates the belief turns out to be a \emph{natural gradient}
step~\cite{Amari98} on the underlying likelihood --- to machine precision, verified over $20{,}000$ random
instances at a worst discrepancy of $1.1 \times 10^{-16}$. Better, it is an
\emph{implicit} step in the sense of stochastic approximation~\cite{RobbinsMonro51},
and therefore unconditionally stable: it cannot leave the
valid range of values for any observation, where the explicit version can. The
confidence coordinate is not a second belief. It is the metric with respect to
which the first one is being learned.

The score that the gradient is assembled from decomposes into contributions
computable at each communication site, from what the site already has --- which
is the identity that makes policy-gradient methods
implementable~\cite{Williams92} and, here, makes the gradient a message rather
than a backward pass. The developer's version of that sentence: grading a
condition gives you gradient descent in the small, correctly conditioned, without importing a learning
framework and without a separate training loop. The learning rule is the
clause.

\subsection{What is settled and what is not}

Three honest boundaries, since the previous subsection was enthusiastic.

\paragraph{The stochastic reading is solid; the quantum reading has a fence
around it.} The complex-valued case is established for a restricted fragment ---
finite, contractive, staged --- and not for the language as a whole. Two
specific hazards sit behind that fence. Left unrestricted, the construction
buys post-selection, and a language with free post-selection decides more than
quantum computers do~\cite{Aaronson05}. And the coherent branching this needs
is quantum \emph{control} rather than the architecture of the quantum process
calculi, which keep the quantum state beside the process rather than in the
branching~\cite{GayNagarajan05,YingYuFeng14,BadescuPanangaden15}. The
research note that develops it~\cite{GWQuantum} states its own boundary as: a finite, contractive,
staged fragment of graded \rhoc{} admits a candidate quantum interpretation, and
the extension to unrestricted reflective concurrency remains open. That fence is
where it is for good reasons, and a developer should read the complex row of the
table above as ``there is a real construction here for circuit-shaped programs''
rather than ``\flang{} is a quantum language''.

\paragraph{Learning under a moving target is not solved.} The measured
experiments include a regime shift, and under it every arm's accuracy degrades
in proportion to how much evidence it had accumulated. Accumulated confidence
makes a learner rigid, and rigidity is exactly wrong when the world changes.
This is a known gap, named as such in the companion note~\cite{GWPLN}, and it is the most
consequential one.

\paragraph{Two searches, not one.} Grading gives a gradient \emph{within} a
hypothesis --- tuning the numbers in a clause. It does not give a gradient
\emph{between} hypotheses --- changing which formula the clause contains. That
second search is discrete and looks like recombination rather than
descent~\cite{ComposeNote}. The
split between the two is not an engineering decision; it follows from the shape
of the hypothesis space.
